\documentclass[11pt]{article}

\usepackage[a4paper,margin=24mm]{geometry}
\usepackage{amsmath,amssymb,amsthm,mathtools}
\usepackage{xcolor}
\usepackage{microtype}
\usepackage[pdfusetitle,colorlinks=true,linkcolor=blue,citecolor=blue,
urlcolor=blue]{hyperref}

\newtheorem{theorem}{Theorem}[section]
\newtheorem{lemma}[theorem]{Lemma}
\newtheorem{proposition}[theorem]{Proposition}
\newtheorem{corollary}[theorem]{Corollary}
\theoremstyle{remark}
\newtheorem{remark}[theorem]{Remark}

\newcommand{\E}{\mathbb E}
\newcommand{\Var}{\operatorname{Var}}
\newcommand{\tr}{\operatorname{tr}}
\newcommand{\I}{\mathbb I}
\newcommand{\N}{\mathcal N}
\newcommand{\calI}{\mathcal I}
\newcommand{\norm}[1]{\left\lVert #1\right\rVert}

\title{A Central Limit Theorem for the High-Temperature SYK
Model}
\author{}
\date{}

\begin{document}
\maketitle

\begin{abstract}
Fix an even interaction order $q\geq4$, and let $Z_N(\beta)$ denote the
partition function of the Gaussian Sachdev--Ye--Kitaev model with $N$
Majorana fermions.  For every sufficiently small fixed inverse temperature
$\beta\neq0$, we prove that
\[
  \frac{Z_N(\beta)-\E Z_N(\beta)}
       {\sqrt{\Var Z_N(\beta)}}
  \Longrightarrow \N(0,1).
\]
We also identify the limiting variance of the relative partition function
and obtain the corresponding central limit theorem for the free energy.  The
leading fluctuation is the radial component of the second Gaussian chaos.
{
The higher chaoses are controlled directly by encoding the sign-gauge parity
rule in the moments of an auxiliary Rademacher chaos.
}
\end{abstract}

\noindent\emph{Keywords:} Sachdev--Ye--Kitaev model,
free energy, central limit theorem, Gaussian chaos.

\medskip
\noindent\emph{Mathematics Subject Classification (2020):}
60F05, 60B20, 82B44.

\section{Introduction and main results}
\label{sec:introduction}

The Sachdev--Ye--Kitaev (SYK) model is a mean-field model of randomly
interacting fermions originating in the work of Sachdev and Ye
\cite{SachdevYe1993} and its Majorana formulation by Kitaev.  Its combination
of solvability, strong interactions, and chaotic behavior has made it a
central model in quantum statistical mechanics and holography; see, for
example, \cite{MaldacenaStanford2016}.  At positive temperature, the
partition function contains both the thermodynamic information and the
sample-to-sample fluctuations of the disorder.

Rigorous results have recently identified the high-temperature annealed and
quenched free-energy limits \cite{GPSZ2026} and established complex
zero-freeness and second-moment estimates \cite{Zlokapa2026}.  The purpose of
this paper is to determine the leading distributional fluctuations in the
same high-temperature regime.  We show that the partition function and its
logarithm have Gaussian fluctuations at scale $N^{1-q/2}$ relative to the
annealed partition function and identify their limiting variance.

\subsection{The SYK model}

Fix an even integer $q\geq4$, and let $N$ be even.  Write
\[
  \calI_N=\{I\subset[N]:|I|=q\},\qquad
  M_N=\binom Nq,\qquad
  \sigma_N^2=\frac{(q-1)!}{N^{q-1}},
\]
and let $(g_I)_{I\in\calI_N}$ be independent standard Gaussian random
variables.  Let $\psi_1,\ldots,\psi_N$ be Majorana matrices satisfying
$\{\psi_i,\psi_j\}=2\delta_{ij}\I$.  For
$I=\{i_1<\cdots<i_q\}$ set
\[
  \psi_I=\mathrm i^{q/2}\psi_{i_1}\cdots\psi_{i_q}.
\]
Then $\psi_I$ is Hermitian, $\psi_I^2=\I$, and $\norm{\psi_I}=1$.  Define
\[
  H_N(g)=\sigma_N\sum_{I\in\calI_N}g_I\psi_I,
  \qquad
  Z_N(z;g)=\tr e^{-zH_N(g)},
  \qquad
  F_N(z)=\E Z_N(z;g).
\]
Here $\tr=2^{-N/2}\operatorname{Tr}$ denotes the normalized trace; in
particular, $Z_N(0)=F_N(0)=1$.  This normalization does not affect the
relative fluctuations studied below.
Throughout, all limits are taken as $N\to\infty$ through even integers.

\subsection{Main results}

\begin{theorem}[Partition-function CLT]
\label{thm:partition-clt}
{
Let $\beta_{\rm Rad}$ be the constant in
Proposition~\ref{prop:higher-chaos-closure}, let $\beta_{\rm FE}$ be the
constant in Proposition~\ref{prop:annealed-free-energy}, and let $\ell_q$ be
the limiting annealed free energy in that proposition.
}
If $\beta\in\mathbb R$ satisfies
\begin{equation}
  {
  0<|\beta|<\min\{\beta_{\rm FE},\beta_{\rm Rad}\},
  \qquad
  \beta\ell_q'(\beta)\neq0.
  }
\end{equation}
Then
\begin{equation}\label{eq:scaled-relative-clt}
  N^{q/2-1}\left(
    \frac{Z_N(\beta)}{\E Z_N(\beta)}-1
  \right)
  \Longrightarrow
  \sqrt{\frac{q!}{2}}\,\beta\ell_q'(\beta)\,\xi,
  \qquad \xi\sim\N(0,1).
\end{equation}
Moreover,
\begin{equation}\label{eq:variance-asymptotic}
  N^{q-2}
  \frac{\Var Z_N(\beta)}{(\E Z_N(\beta))^2}
  \longrightarrow
  \frac{q!}{2}\,\beta^2\ell_q'(\beta)^2,
\end{equation}
and hence
\begin{equation}\label{eq:standardized-clt}
  \frac{Z_N(\beta)-\E Z_N(\beta)}
  {\sqrt{\Var Z_N(\beta)}}
  \Longrightarrow\N(0,1).
\end{equation}
\end{theorem}

\begin{corollary}[Log-partition-function CLT]
\label{cor:free-energy-clt}
Under the assumptions of Theorem~\ref{thm:partition-clt},
\[
  N^{q/2-1}\left(
    \log Z_N(\beta)-\E\log Z_N(\beta)
  \right)
  \Longrightarrow
  \sqrt{\frac{q!}{2}}\,\beta\ell_q'(\beta)\,\xi,
  \qquad \xi\sim\N(0,1).
\]
\end{corollary}

\subsection{Relation to earlier work}

Feng, Tian, and Wei \cite{FTW2018} proved a central limit theorem for linear
eigenvalue statistics of the Gaussian SYK Hamiltonian.  Their theorem treats
a fixed test function $f$ under regularity assumptions on $f$ and $f'$.
In the normalization used in that paper, however, the partition function
corresponds to
\[
  f_{N,\beta}(x)=e^{-\beta\tau_Nx},
  \qquad \tau_N\asymp\sqrt N,
\]
so both the test function and its derivative vary with $N$.  The
fixed-test-function theorem therefore does not apply directly to
$Z_N(\beta)$.

{
Our argument retains the radial-fluctuation mechanism identified in
\cite{FTW2018}, but establishes the uniform remainder estimate needed for
this $N$-dependent sequence.  The sign-gauge rule says that a Hermite
coefficient can be nonzero only when every vertex has even total incidence.
Introducing one independent Rademacher sign per vertex turns the sum over
such multi-indices into the exponential generating function of a single
degree-$q$ Rademacher chaos.  A uniform small-parameter exponential-moment
bound then controls all Gaussian chaoses of order at least three at once.
The proof uses only annealed free-energy convergence and does not assume any
fluctuation theorem for either $Z_N$ or $\log Z_N$.
}

\subsection{Outline of the proof}


Clifford sign-gauge symmetry eliminates the first Gaussian chaos and
identifies the second chaos exactly:
\[
  Y_N(\beta)-1
  =
  A_N(\beta)Q_N+\mathcal R_N(\beta),
  \]
  where
  \[
    A_N(\beta)=\frac{\beta}{2}\partial_\beta\log F_N(\beta),
    \qquad
    Q_N=\frac1{M_N}\sum_I(g_I^2-1).
  \]
Here $\mathcal R_N$ is the sum of the Gaussian chaoses of orders at least
three.  Thus the leading random term is a normalized sum of $M_N$
independent centered variables.  The main task is to prove
\[
  \norm{\mathcal R_N(\beta)}_2^2=o(N^{-(q-2)}).
\]

{
For the higher-order part, a Duhamel derivative bound and the sign-gauge
parity rule give
\begin{equation}
  \sum_{r\geq3}\lVert P_rY_N(\beta)\rVert_2^2
  \leq
  \sum_{r\geq3}t_N^r
  \sum_{\substack{|\alpha|=r\\ \alpha\ \mathrm{even}}}
  \frac{1}{\alpha!},
  \qquad
  t_N=\beta^2\sigma_N^2.
\end{equation}
If $(\varepsilon_x)_{x=1}^N$ are independent Rademacher variables and
$S_{N,q}=\sum_I\prod_{x\in I}\varepsilon_x$, then the inner sum is exactly
$\E_\varepsilon[S_{N,q}^r]/r!$.  Newton's identities give
$|S_{N,q}|\leq C_q(N^{q/2}+|\sum_x\varepsilon_x|^q)$, from which a
sub-Gaussian estimate yields
\begin{equation}
  \E_\varepsilon\!\left[
    |S_{N,q}|^3e^{t_N|S_{N,q}|}
  \right]
  \leq C_qN^{3q/2}
\end{equation}
for sufficiently small fixed $|\beta|$.  The entire higher-order sum is
therefore $O(N^{-3(q/2-1)})=o(N^{-(q-2)})$.
}

The annealed free-energy input and convexity imply
\[
  \frac{A_N(\beta)}N
  \longrightarrow \frac{\beta}{2}\ell_q'(\beta).
\]
Since $M_N\sim N^q/q!$, the classical central limit theorem for $Q_N$, the
remainder estimate  gives the claimed limit for
$N^a(Y_N-1)$.  Orthogonality of distinct Gaussian chaoses gives the matching
variance asymptotic and hence the standardized central limit theorem.

{
\section{Free-energy input}
\label{sec:inputs}
}

\subsection{Annealed free energy and coefficient asymptotics}

To identify the deterministic coefficient $A_N(\beta)$ in
\eqref{eq:finite-N-FTW-coefficient}, we use the annealed free-energy limit
from \cite{GPSZ2026}.  The required convergence of derivatives follows from
convexity and is not an additional input.

\begin{lemma}[Convex convergence and derivatives]
\label{lem:convex-derivatives}
Let $I\subset\mathbb R$ be open.  Suppose that $h_N:I\to\mathbb R$ are finite,
differentiable, convex functions and $h_N(x)\to h(x)$ pointwise, where
$h\in C^1(I)$.  Then $h_N\to h$ and $h_N'\to h'$ locally uniformly on $I$.
\end{lemma}

\begin{proof}
Fix a compact interval $[a,b]\subset I$ and choose $u<a<b<v$ in $I$.
Monotonicity of secant slopes bounds every secant slope of $h_N$ on $[a,b]$
between quantities formed from $h_N(u),h_N(a),h_N(b),h_N(v)$.  These
endpoint values converge, so the functions $h_N$ are equi-Lipschitz on
$[a,b]$.  Pointwise convergence on a finite mesh, followed by mesh
refinement, gives uniform convergence there.

For $x\in I$ and sufficiently small $s>0$, convexity gives
\[
  \frac{h_N(x)-h_N(x-s)}s
  \leq h_N'(x)\leq
  \frac{h_N(x+s)-h_N(x)}s.
\]
First let $N\to\infty$ and then $s\downarrow0$ to obtain
$h_N'(x)\to h'(x)$.  Finally, $h_N'$ are nondecreasing.  On a fine partition
of $[a,b]$, monotonicity traps $h_N'(x)$ between its values at adjacent mesh
points, while uniform continuity traps $h'(x)$ between the corresponding
values of $h'$.  Pointwise convergence at the finitely many mesh points then
gives uniform convergence of the derivatives.
\end{proof}

\begin{proposition}[High-temperature annealed free energy]
\label{prop:annealed-free-energy}
There are $\beta_{\rm FE}=\beta_{\rm FE}(q)>0$ and an even real-analytic
function $\ell_q$ on $(-\beta_{\rm FE},\beta_{\rm FE})$ such that, locally
uniformly on this interval,
\begin{align}
  \frac1N\log F_N(\beta)&\longrightarrow\ell_q(\beta),
  \label{eq:annealed-free-energy-limit}\\
  \frac1N\partial_\beta\log F_N(\beta)&\longrightarrow\ell_q'(\beta).
  \label{eq:annealed-derivative-limit}
\end{align}
Moreover,
\begin{equation}\label{eq:ell-small-beta}
  \ell_q(\beta)=\frac{\beta^2}{2q}+O(\beta^4),
  \qquad
  \ell_q'(\beta)=\frac\beta q+O(\beta^3).
\end{equation}
\end{proposition}

\begin{proof}
The pointwise convergence in \eqref{eq:annealed-free-energy-limit} is Theorem
5 of \cite{GPSZ2026}, together with the Gaussian extension stated there.  That
extension invokes the Gaussian--Rademacher universality result
\cite{AGK2025}.

We first record the normalization conversion.  The
Majoranas of \cite{GPSZ2026} satisfy
\[
  \psi_i^{\rm GPSZ}=2^{-1/2}\psi_i,
  \qquad
  H_N^{\rm GPSZ}=2^{-q/2}H_N.
\]
Set $u=2^{q/2}\beta$.  Gaussian symmetry gives
\[
  \E\tr e^{-\beta H_N}
  =
  \E\tr e^{+\beta H_N},
\]
and $uH_N^{\rm GPSZ}=\beta H_N$.  Thus the sign and normalization conventions
of \cite{GPSZ2026} give
\begin{equation}\label{eq:GPSZ-conversion}
  \frac1N\log F_N(\beta)=\frac{u}{N}f_{N,q}^{\rm GPSZ}(u),
\end{equation}
where $f_{N,q}^{\rm GPSZ}(u)=u^{-1}\log\E Z_N^{\rm GPSZ}(u)$.

We next extract the regularity and the first coefficient from the explicit
formula in that theorem.  In its notation,
\[
  \Phi_u(z)=\sum_{v\geq1}\mu_v(u)z^v,
  \qquad
  u f_q^{\rm GPSZ}(u)=\Phi_u(t_*(u))-1-\log t_*(u),
\]
where $t_*\partial_z\Phi_u(t_*)=1$.  A contributing hypertree has
$v=1+(q-1)s$ vertices and its coefficient contains
\[
  \rho(u)^{v-1},
  \qquad
  \rho(u)^{q-1}=\frac{(q-1)!u^2}{2^{q+1}}.
\]
More explicitly, $\mathcal T_v^{\rm sim}$ is empty unless
$v=1+(q-1)s$ for some $s\geq0$.  For $s\geq1$ we may therefore write
\[
  \mu_{1+(q-1)s}(u)=c_{q,s}u^{2s},
\]
where
\[
  c_{q,s}
  =
  \frac1{(1+(q-1)s)!}
  \left(\frac{(q-1)!}{2^{q+1}}\right)^s
  \sum_{T\in\mathcal T_{1+(q-1)s}^{\rm sim}}b_T.
\]
The bounds used in the proof of Proposition 3 of \cite{GPSZ2026} give
$|c_{q,s}|\leq C_q^s$.  Hence
\[
  \Phi_u(z)
  =
  z+\sum_{s\geq1}
  c_{q,s}u^{2s}z^{1+(q-1)s}
\]
converges normally for $(u^2,z)$ in a neighborhood of $(0,1)$.  Thus
$\Phi_u(z)$ is jointly analytic in $(u^2,z)$ there.  At
$u=0$, $\Phi_0(z)=z$ and $t_*(0)=1$; the derivative with respect to $t$ of
$t\partial_z\Phi_u(t)-1$ is $1$ at $(u,t)=(0,1)$.  The analytic implicit
function theorem makes $t_*$, and hence the limiting right-hand side, analytic
in $u^2$.

The first nontrivial term is the unique one-edge hypertree on $q$ vertices:
\[
  \Phi_u(z)=z+\frac{u^2}{q2^{q+1}}z^q+O(u^4).
\]
It follows successively that
\[
  t_*(u)=1-\frac{u^2}{2^{q+1}}+O(u^4),
  \qquad
  u f_q^{\rm GPSZ}(u)=\frac{u^2}{q2^{q+1}}+O(u^4).
\]
Using $u=2^{q/2}\beta$ in \eqref{eq:GPSZ-conversion} gives
\eqref{eq:ell-small-beta}.

It remains only to justify the two local-uniform assertions.  Let
$h_N=N^{-1}\log F_N$.  If
\[
  \overline\nu_N(B)=\E\,\tr\mathbf1_{\{H_N\in B\}},
\]
then $F_N(\beta)=\int e^{-\beta x}\overline\nu_N(\mathrm dx)$; hence $h_N$
is finite, differentiable, and convex.  {Differentiation is
justified by
$\norm{H_N(g)}\leq\sigma_N\sum_{I\in\calI_N}|g_I|$ and Gaussian
exponential integrability.}
The limiting function just obtained is $C^1$.  Lemma
\ref{lem:convex-derivatives} now proves both local-uniform convergences.
\end{proof}

By \eqref{eq:finite-N-FTW-coefficient} and
\eqref{eq:annealed-derivative-limit},
\begin{equation}\label{eq:AN-asymptotic}
  \frac{A_N(\beta)}N\longrightarrow
  \frac\beta2\ell_q'(\beta).
\end{equation}

\section{Gaussian-chaos decomposition}
\label{sec:chaos-decomposition}
Set
\begin{equation}
  a=\frac q2-1>0,
  \qquad
  Y_N(\beta;g)=\frac{Z_N(\beta;g)}{F_N(\beta)}.
\end{equation}
Thus $\E Y_N(\beta)=1$.  Let $P_r$ denote the orthogonal projection in
$L^2(\gamma_N)$ onto the $r$-th homogeneous Gaussian chaos, where $\gamma_N$
denotes standard Gaussian measure on $\mathbb R^{M_N}$.  We write
$\norm{X}_2^2=\E|X|^2$.  We use the probabilists' Hermite polynomials
\begin{equation}
  H_\alpha(g)=\prod_{I\in\calI_N}{\rm He}_{\alpha_I}(g_I),
  \qquad
  \E[H_\alpha H_\eta]
  =\alpha!\,\mathbf 1_{\{\alpha=\eta\}},
\end{equation}
where $\alpha!=\prod_I\alpha_I!$.  For every sufficiently regular random
variable $X$,
\begin{equation}\label{eq:Hermite-coefficient-formula}
  \norm{P_rX}_2^2
  =
  \sum_{|\alpha|=r}
  \frac{|\E\,\partial_g^\alpha X|^2}{\alpha!}.
\end{equation}

We first relate our normalization to that of Feng, Tian, and Wei.  Define
\begin{equation}\label{eq:FTW-normalization}
  \widehat H_N(g)=M_N^{-1/2}\sum_{I\in\calI_N}g_I\psi_I,
  \qquad
  \tau_N=\sigma_N\sqrt{M_N}.
\end{equation}
Then $H_N=\tau_N\widehat H_N$ and
\begin{equation}\label{eq:tau-asymptotic}
  \tau_N^2=\sigma_N^2M_N=\frac Nq\bigl(1+O_q(N^{-1})\bigr).
\end{equation}
Consequently, for
$\widehat L_N(f)=\tr f(\widehat H_N)$, then
\begin{equation}\label{eq:partition-as-linear-statistic}
  Z_N(\beta)=\widehat L_N(f_{N,\beta}),
  \qquad
  f_{N,\beta}(x)=e^{-\beta\tau_Nx}.
\end{equation}

For a fixed test function $f$, Feng--Tian--Wei \cite{FTW2018} prove the
Gaussian SYK linear-statistics CLT by isolating
\begin{equation}\label{eq:FTW-radial-statistic}
  Q_N:=\tr\widehat H_N^2-\E\tr\widehat H_N^2
      =\frac1{M_N}\sum_{I\in\calI_N}(g_I^2-1).
\end{equation}
In the fixed-$q$ regime, their coefficient is
\[
  A(f)=\int_{\mathbb R}\frac{x f'(x)}2\,\rho_\infty(\mathrm dx),
\]
where $\rho_\infty$ is the limiting spectral measure.
They show that the centered statistic associated with
$f-A(f)x^2$ has variance $o(M_N^{-1})$.  Since
$\sqrt{M_N}Q_N\Rightarrow\N(0,2)$, the radial term gives their limit.

There are two obstructions to substituting \eqref{eq:partition-as-linear-statistic}
directly into that theorem.  First, $f_{N,\beta}$ depends on $N$ through
$\tau_N\asymp\sqrt N$.  Second, neither the fixed-function polynomial
approximation nor the bounded-derivative estimate in \cite{FTW2018} is uniform
for this sequence.  In particular, the theorem as stated in
\cite{FTW2018} does not imply a CLT for $Z_N(\beta)$.

We instead work with the relative partition function $Y_N$ and replace the
limiting coefficient $A(f)$ by the exact finite-$N$ coefficient
\begin{equation}\label{eq:finite-N-FTW-coefficient}
  A_N(\beta):=\frac{\beta}{2}\,\partial_\beta\log F_N(\beta).
\end{equation}
The next lemma proves the exact radial decomposition
\begin{equation}\label{eq:exact-FTW-decomposition-preview}
  Y_N(\beta)-1=A_N(\beta)Q_N+\mathcal R_N(\beta),
  \qquad
  \mathcal R_N=\sum_{r\geq3}P_rY_N.
\end{equation}
The estimate replacing the fixed-test-function variance computation of
\cite{FTW2018} is
\begin{equation}\label{eq:target-FTW-remainder}
  \E|\mathcal R_N(\beta)|^2=o(N^2/M_N)=o(N^{2-q}).
\end{equation}
{
Subsection~\ref{sec:Rademacher-remainder} proves
\eqref{eq:target-FTW-remainder} directly.
}

\subsection{Exact finite-size radial projection}

\begin{lemma}[First and second chaoses]\label{lem:exact-second-chaos}
For every real $\beta$,
\[
  P_1Y_N(\beta)=0
\]
and
\begin{equation}\label{eq:P2Y-exact}
  P_2Y_N(\beta)
  =
  \frac{\beta\,\partial_\beta\log F_N(\beta)}{2M_N}
  \sum_{I\in\calI_N}(g_I^2-1).
\end{equation}
Equivalently, Clifford trace orthogonality gives
\[
  \tr(\psi_I\psi_J)=\mathbf 1_{\{I=J\}},
  \qquad
  \tr H_N^2=\sigma_N^2\sum_Ig_I^2.
\]
so the second chaos is the $L^2$ projection of $Y_N-1$ onto the centered
radial statistic
$\tr H_N^2-\E\tr H_N^2$.
\end{lemma}

\begin{proof}
For $s=(s_1,\ldots,s_N)\in\{\pm1\}^N$, define
$\chi_s(I)=\prod_{x\in I}s_x$.  We first prove in detail that
\begin{equation}\label{eq:sign-gauge-invariance}
  Z_N(\beta;g)=Z_N\bigl(\beta;(\chi_s(I)g_I)_I\bigr).
\end{equation}
To make the Clifford symmetry explicit, let
\[
  \Gamma
  := 
  \psi_1\cdots\psi_N.
\]
Note that $\Gamma^*\Gamma$ is identity map, so $\Gamma$ is unitary. Since $N$ is even, moving $\psi_j$ through the other $N-1$ Majorana
generators produces the sign $(-1)^{N-1}=-1$.  Hence
\[
  \Gamma\psi_j\Gamma^*=-\psi_j
  \qquad (j\in[N]).
\]
For each $x\in[N]$, put $U_x=\Gamma\psi_x$.  This is unitary, and conjugation
by $U_x$ is the composition of conjugation by $\psi_x$ and conjugation by
$\Gamma$.  The Clifford relations give
\[
  \psi_x\psi_j\psi_x
  =
  \begin{cases}
    \psi_x,&j=x,\\
    -\psi_j,&j\neq x.
  \end{cases}
\]
Combining the last two displays, we obtain
\[
  U_x\psi_jU_x^*
  =
  \begin{cases}
    -\psi_x,&j=x,\\
    \psi_j,&j\neq x.
  \end{cases}
\]
Thus conjugation by $U_x$ changes the sign of precisely the generator
$\psi_x$.

Let $S_s=\{x\in[N]:s_x=-1\}$ and let $U_s$ be the product of the unitaries
$U_x$ over $x\in S_s$, in any fixed order.  Successively applying the
preceding conjugation rule yields
\[
  U_s\psi_jU_s^*=s_j\psi_j
  \qquad (j\in[N]).
\]
If $I=\{i_1<\cdots<i_q\}$, then conjugation is multiplicative, so the
definition $\psi_I=\mathrm i^{q/2}\psi_{i_1}\cdots\psi_{i_q}$ gives
\begin{align*}
  U_s\psi_IU_s^*
  &=
  \mathrm i^{q/2}
  (U_s\psi_{i_1}U_s^*)\cdots
  (U_s\psi_{i_q}U_s^*)\\
  &=
  \left(\prod_{x\in I}s_x\right)\psi_I
  =\chi_s(I)\psi_I.
\end{align*}
Writing $g^{(s)}=(\chi_s(I)g_I)_I$, we therefore have
\begin{align*}
  H_N(g^{(s)})
  &=\sigma_N\sum_{I\in\calI_N}\chi_s(I)g_I\psi_I\\
  &=\sigma_N\sum_{I\in\calI_N}g_IU_s\psi_IU_s^*
    =U_sH_N(g)U_s^*.
\end{align*}
The power-series definition of the exponential then implies
\[
  e^{-\beta H_N(g^{(s)})}
  =U_se^{-\beta H_N(g)}U_s^*.
\]
Finally, the normalized trace is invariant under unitary conjugation.
Consequently
\[
  Z_N(\beta;g^{(s)})
  =\tr\!\left(U_se^{-\beta H_N(g)}U_s^*\right)
  =\tr e^{-\beta H_N(g)}
  =Z_N(\beta;g),
\]
which proves \eqref{eq:sign-gauge-invariance}.

More precisely,
\[
  H_\alpha((\chi_s(I)g_I)_I)
  =
  \left(
    \prod_{x=1}^N
    s_x^{\sum_{I\ni x}\alpha_I}
  \right)H_\alpha(g).
\]
Averaging the Hermite coefficient against all $s$ shows that it vanishes
unless $\sum_{I\ni x}\alpha_I$ is even for every $x$.  This parity rule kills
the first chaos.  In the second chaos it also kills every off-diagonal
monomial $g_Ig_J$, since the parity condition for such a monomial forces
$I=J$.
Permutation symmetry shows that the remaining diagonal coefficients agree.
Consequently
\[
  P_2Z_N(\beta)=b_N(\beta)\sum_I(g_I^2-1)
\]
for a deterministic scalar $b_N$.

It remains to compute the coefficient.  For each fixed $I$, the
probabilists' second Hermite polynomial is $g_I^2-1$, with
\[
  \E(g_I^2-1)^2=2.
\]
Since $P_2$ is an orthogonal projection and the variables $g_I$ are
independent, taking the inner product of the preceding formula for $P_2Z_N$
with $g_I^2-1$ gives
\[
  2b_N(\beta)
  =\E\bigl[Z_N(\beta)(g_I^2-1)\bigr].
\]
Applying Gaussian integration by parts twice, conditionally on all
coordinates other than $g_I$, yields
\[
  \E\bigl[Z_N(\beta)(g_I^2-1)\bigr]
  =\E\,\partial_{g_I}^2Z_N(\beta).
\]
Summing over the $M_N$ choices of $I$ therefore gives the first equality
below.  A further conditional Gaussian integration by parts gives the second:
\begin{align*}
  2M_Nb_N(\beta)
  &=\sum_I\E\,\partial_{g_I}^2Z_N(\beta)
    =\sum_I\E\bigl[g_I\partial_{g_I}Z_N(\beta)\bigr]\\
  &=\beta\,\partial_\beta\E Z_N(\beta)
    =\beta F_N'(\beta).
\end{align*}

The last equality follows pointwise from the exact scaling identity
\[
  Z_N(\beta;tg)
  =\tr e^{-\beta H_N(tg)}
  =\tr e^{-t\beta H_N(g)}
  =Z_N(t\beta;g),
\]
where we used $H_N(tg)=tH_N(g)$.  Differentiating at $t=1$ gives
\[
  \sum_Ig_I\partial_{g_I}Z_N(\beta;g)
  =\beta\,\partial_\beta Z_N(\beta;g).
\]
Finally,
\[
  \partial_\beta Z_N(\beta;g)
  =-\tr\!\left(H_N(g)e^{-\beta H_N(g)}\right),
  \qquad
  \left|\partial_\beta Z_N(\beta;g)\right|
  \leq \norm{H_N(g)}e^{|\beta|\norm{H_N(g)}}.
\]
Together with
\[
  \norm{H_N(g)}\leq\sigma_N\sum_I|g_I|,
\]
the right-hand side has every exponential moment under the finite dimensional
Gaussian law.  It gives a locally uniform dominating function in $\beta$.
The same type of bound justifies the preceding Gaussian integrations by
parts.  Thus differentiation passes through the expectation and
$\partial_\beta\E Z_N(\beta)=F_N'(\beta)$.
Dividing by the deterministic quantity $F_N(\beta)$ proves
\eqref{eq:P2Y-exact}.
\end{proof}

{
\subsection{Direct Rademacher control of all higher chaoses}
\label{sec:Rademacher-remainder}
}

\begin{lemma}[Real-temperature derivative bound]
\label{lem:Z-derivative-bound}
For real $\beta$, every multi-index $\alpha$ of total order $r$ satisfies
\begin{equation}\label{eq:Z-derivative-bound}
  \left|\partial_g^\alpha Z_N(\beta;g)\right|
  \leq (|\beta|\sigma_N)^r Z_N(\beta;g).
\end{equation}
\end{lemma}

\begin{proof}
Label the $r$ differentiations by $I_1,\ldots,I_r$.  The iterated Duhamel
formula is a sum over $r!$ permutations of integrals over the simplex
$\Delta_r=\{t_j\geq0:\sum_{j=0}^rt_j=1\}$, whose volume is $1/r!$.
Every integrand, apart from the factor $(|\beta|\sigma_N)^r$, has the form
\[
  \tr\left(
    e^{-t_0\beta H_N}\psi_{I_{\pi(1)}}e^{-t_1\beta H_N}
    \cdots
    \psi_{I_{\pi(r)}}e^{-t_r\beta H_N}
  \right).
\]
Here and below the Schatten norms are defined using the normalized trace:
\[
  \norm{A}_p=(\tr|A|^p)^{1/p},
  \qquad
  \norm{A}_\infty=\norm{A}.
\]
For $t_j>0$,
\[
  \norm{e^{-t_j\beta H_N}}_{1/t_j}
  =
  \bigl(\tr e^{-\beta H_N}\bigr)^{t_j},
\]
with the usual $p=\infty$ interpretation when $t_j=0$.  Noncommutative
H\"older, with Schatten exponents $1/t_j$, and
$\norm{\psi_I}_\infty=1$, therefore bounds the absolute value of the
integrand by
\[
  \prod_{j=0}^r
  \left(\tr e^{-\beta H_N}\right)^{t_j}
  =Z_N(\beta;g).
\]
The permutation count cancels the simplex volume, proving
\eqref{eq:Z-derivative-bound}.  The same labeled formula covers repeated
coordinates.
\end{proof}

\begingroup


Call a multi-index $\alpha$ \emph{even} if
\begin{equation}\label{eq:even-multi-index}
  \sum_{I\ni x}\alpha_I\equiv0\pmod2
  \qquad\text{for every }x\in[N].
\end{equation}
The sign-gauge symmetry used in Lemma~\ref{lem:exact-second-chaos} shows that
only even multi-indices have nonzero Hermite coefficients.

Let $(\varepsilon_x)_{x=1}^N$ be independent Rademacher variables on an
auxiliary probability space, and write $\E_\varepsilon$ for expectation on
this space.  Set
\begin{equation}\label{eq:Rademacher-chaos-definition}
  L_N=\sum_{x=1}^N\varepsilon_x,
  \qquad
  S_{N,q}
  =\sum_{I\in\calI_N}\varepsilon_I,
  \qquad
  \varepsilon_I=\prod_{x\in I}\varepsilon_x.
\end{equation}
Thus $S_{N,q}$ is the degree-$q$ elementary symmetric polynomial in the
Rademacher signs.

\begin{lemma}[Rademacher encoding and exponential moment]
\label{lem:Rademacher-encoding}
For every integer $r\geq0$,
\begin{equation}\label{eq:even-sum-Rademacher-moment}
  \sum_{\substack{|\alpha|=r\\ \alpha\ \mathrm{even}}}
  \frac1{\alpha!}
  =\frac1{r!}\E_\varepsilon\!\left[S_{N,q}^r\right].
\end{equation}
Moreover, there are constants $c_q>0$ and $C_q<\infty$, depending only on
$q$, such that
\begin{equation}\label{eq:Rademacher-exponential-moment}
  \E_\varepsilon\!\left[
    |S_{N,q}|^3e^{t|S_{N,q}|}
  \right]
  \leq C_qN^{3q/2}
\end{equation}
for every $N$ and every $0\leq t\leq c_qN^{-(q-1)}$.
\end{lemma}

\begin{proof}
Expanding $S_{N,q}^r$ by the multinomial formula gives
\begin{equation}\label{eq:Rademacher-multinomial-expansion}
  \frac1{r!}\E_\varepsilon\!\left[S_{N,q}^r\right]
  =
  \sum_{|\alpha|=r}\frac1{\alpha!}
  \E_\varepsilon\!\left[
    \prod_{I\in\calI_N}\varepsilon_I^{\alpha_I}
  \right].
\end{equation}
For a fixed $\alpha$, the product inside the expectation is
\begin{equation}
  \prod_{x=1}^N
  \varepsilon_x^{\sum_{I\ni x}\alpha_I}.
\end{equation}
Its expectation is one if $\alpha$ is even and zero otherwise.  This proves
\eqref{eq:even-sum-Rademacher-moment}.

We next prove a deterministic bound for $S_{N,q}$.  For $0\leq m\leq q$,
let
\begin{equation}
  e_m(\varepsilon)
  =\sum_{\substack{I\subset[N]\\ |I|=m}}
  \prod_{x\in I}\varepsilon_x,
  \qquad e_0(\varepsilon)=1.
\end{equation}
The power sums of the Rademacher signs satisfy
\begin{equation}
  p_j:=\sum_{x=1}^N\varepsilon_x^j
  =
  \begin{cases}
    L_N,&j\text{ odd},\\
    N,&j\text{ even}.
  \end{cases}
\end{equation}
Newton's identities state that
\begin{equation}\label{eq:Newton-identities}
  m e_m
  =\sum_{j=1}^m(-1)^{j-1}e_{m-j}p_j.
\end{equation}
An induction on $m$, using \eqref{eq:Newton-identities}, shows that there is
a constant $C_m<\infty$ such that
\begin{equation}\label{eq:elementary-symmetric-bound}
  |e_m(\varepsilon)|
  \leq C_m\left(N^{m/2}+|L_N|^m\right).
\end{equation}
For completeness, suppose that the estimate is known below order $m$ and
write $u=|L_N|/\sqrt N$.  If $j$ is even, then
\begin{align}
  |e_{m-j}p_j|
  &\leq
  C_{m-j}\left(N^{(m-j)/2}+|L_N|^{m-j}\right)N \\
  &\leq
  C_{m-j}N^{m/2}\left(1+u^m\right),
\end{align}
because $j\geq2$ and $u^{m-j}\leq1+u^m$.  If $j$ is odd, then
\begin{align}
  |e_{m-j}p_j|
  &\leq
  C_{m-j}\left(N^{(m-j)/2}+|L_N|^{m-j}\right)|L_N| \\
  &\leq
  C_{m-j}N^{m/2}\left(1+u^m\right),
\end{align}
because $j\geq1$ and $u^{m-j+1}\leq1+u^m$.  Substitution into
\eqref{eq:Newton-identities} proves the induction step.  Since
$S_{N,q}=e_q(\varepsilon)$, we have in particular
\begin{equation}\label{eq:Rademacher-chaos-deterministic-bound}
  |S_{N,q}|
  \leq C_q\left(N^{q/2}+|L_N|^q\right).
\end{equation}

Set $X_N=L_N/\sqrt N$.  Since $|X_N|\leq\sqrt N$ and
$a=q/2-1$, \eqref{eq:Rademacher-chaos-deterministic-bound} implies that,
for $0\leq t\leq c_qN^{-(q-1)}$,
\begin{align}
  t|S_{N,q}|
  &\leq C_qc_qN^{-a}\left(1+|X_N|^q\right) \\
  &\leq C_qc_q\left(1+|X_N|^2\right).
  \label{eq:Rademacher-exponent-reduction}
\end{align}
Indeed,
\begin{equation}
  N^{-a}|X_N|^q
  =|X_N|^2
  \left(\frac{|X_N|}{\sqrt N}\right)^{q-2}
  \leq |X_N|^2.
\end{equation}
Choose $c_q$ small enough that the last member of
\eqref{eq:Rademacher-exponent-reduction} is bounded by
$(1+|X_N|^2)/4$.

Hoeffding's inequality gives
\begin{equation}\label{eq:Hoeffding-Rademacher-sum}
  \mathbb P_\varepsilon(|X_N|\geq u)
  \leq2e^{-u^2/2}
  \qquad (u\geq0).
\end{equation}
Consequently, integration of the tail in
\eqref{eq:Hoeffding-Rademacher-sum} yields
\begin{equation}\label{eq:uniform-Rademacher-integrability}
  \sup_{N\geq1}
  \E_\varepsilon\!\left[
    \left(1+|X_N|^{3q}\right)e^{|X_N|^2/4}
  \right]
  <\infty.
\end{equation}
Finally, \eqref{eq:Rademacher-chaos-deterministic-bound},
\eqref{eq:Rademacher-exponent-reduction}, and
\eqref{eq:uniform-Rademacher-integrability} give
\begin{align}
  \E_\varepsilon\!\left[
    |S_{N,q}|^3e^{t|S_{N,q}|}
  \right]
  &\leq
  C_qN^{3q/2}
  \E_\varepsilon\!\left[
    \left(1+|X_N|^{3q}\right)
    e^{(1+|X_N|^2)/4}
  \right] \\
  &\leq C_qN^{3q/2},
\end{align}
which proves \eqref{eq:Rademacher-exponential-moment}.
\end{proof}

\begin{proposition}[Higher-chaos closure]
\label{prop:higher-chaos-closure}
There exists $\beta_{\rm Rad}=\beta_{\rm Rad}(q)>0$ such that, for every
fixed real $\beta$ with $|\beta|<\beta_{\rm Rad}$,
\begin{equation}\label{eq:higher-chaos-closure}
  \sum_{r\geq3}\norm{P_rY_N(\beta)}_2^2
  =O\!\left(N^{-3(q/2-1)}\right)
  =o(N^{2-q}).
\end{equation}
\end{proposition}

\begin{proof}
Put
\begin{equation}\label{eq:tN-definition}
  t_N=\beta^2\sigma_N^2
  =\beta^2(q-1)!N^{-(q-1)}.
\end{equation}
The Hermite coefficient formula, Lemma~\ref{lem:Z-derivative-bound}, and
$F_N(\beta)>0$ give
\begin{align}
  \norm{P_rY_N(\beta)}_2^2
  &=\frac1{F_N(\beta)^2}
    \sum_{|\alpha|=r}\frac1{\alpha!}
    \left|\E\partial_g^\alpha Z_N(\beta)\right|^2 \\
  &\leq t_N^r
    \sum_{\substack{|\alpha|=r\\ \alpha\ \mathrm{even}}}
    \frac1{\alpha!}.
  \label{eq:Hermite-Rademacher-bound}
\end{align}
Here the sum may be restricted to even $\alpha$ by the sign-gauge symmetry.
By \eqref{eq:even-sum-Rademacher-moment}, summing
\eqref{eq:Hermite-Rademacher-bound} over all $r\geq3$ yields
\begin{align}
  \sum_{r\geq3}\norm{P_rY_N(\beta)}_2^2
  &\leq
  \sum_{r\geq3}\frac{t_N^r}{r!}
  \E_\varepsilon\!\left[S_{N,q}^r\right] \\
  &\leq
  \E_\varepsilon\!\left[
    \sum_{r\geq3}\frac{(t_N|S_{N,q}|)^r}{r!}
  \right] \\
  &\leq
  \frac{t_N^3}{6}
  \E_\varepsilon\!\left[
    |S_{N,q}|^3e^{t_N|S_{N,q}|}
  \right].
  \label{eq:all-higher-chaoses-at-once}
\end{align}
Choose
\begin{equation}
  \beta_{\rm Rad}
  =\sqrt{\frac{c_q}{(q-1)!}},
\end{equation}
where $c_q$ is the constant in Lemma~\ref{lem:Rademacher-encoding}.  If
$|\beta|<\beta_{\rm Rad}$, then
$t_N\leq c_qN^{-(q-1)}$.  Therefore
\eqref{eq:Rademacher-exponential-moment} and
\eqref{eq:all-higher-chaoses-at-once} imply
\begin{align}
  \sum_{r\geq3}\norm{P_rY_N(\beta)}_2^2
  &\leq C_{q,\beta}
  N^{-3(q-1)}N^{3q/2} \\
  &=C_{q,\beta}N^{-3(q/2-1)}.
\end{align}
Since $q\geq4$, we have $3(q/2-1)>q-2$, which completes the proof.
\end{proof}

\endgroup
\section{Proof of the main results}
\begin{proof}[Proof of Theorem~\ref{thm:partition-clt}]
Let
\[
  S_N=\frac1{\sqrt{2M_N}}\sum_{I\in\calI_N}(g_I^2-1).
\]
The classical i.i.d.\ central limit theorem gives
$S_N\Rightarrow\xi\sim\N(0,1)$.
By Lemma~\ref{lem:exact-second-chaos},
\begin{equation}\label{eq:Y-decomposition-final}
  Y_N(\beta)-1
  =c_N(\beta)S_N+\mathcal R_N(\beta),
  \qquad
  c_N(\beta)
  =\frac{\beta\,\partial_\beta\log F_N(\beta)}{\sqrt{2M_N}},
\end{equation}
where $\mathcal R_N=\sum_{r\geq3}P_rY_N$ is orthogonal to $S_N$.
Since $M_N\sim N^q/q!$, \eqref{eq:annealed-derivative-limit} gives
\[
  N^{q/2-1}c_N(\beta)
  \longrightarrow
  \sqrt{\frac{q!}{2}}\,\beta\ell_q'(\beta).
\]
Proposition~\ref{prop:higher-chaos-closure} gives
$N^{q/2-1}\mathcal R_N\to0$ in $L^2$.  This proves
\eqref{eq:scaled-relative-clt}.  Orthogonality in
\eqref{eq:Y-decomposition-final} also gives
\[
  \Var(Y_N)=c_N(\beta)^2+\norm{\mathcal R_N(\beta)}_2^2,
\]
which proves \eqref{eq:variance-asymptotic}.  Finally,
$\Var(Y_N)=\Var(Z_N)/(\E Z_N)^2$, and the assumed nonvanishing of the limiting
coefficient give 
\eqref{eq:standardized-clt}.
\end{proof}

\begin{proof}[Proof of Corollary~\ref{cor:free-energy-clt}]
Put
\[
  a=\frac q2-1,\qquad
  X_N=\frac{Z_N(\beta)-\E Z_N(\beta)}{\E Z_N(\beta)},
  \qquad
  W_N=\log\frac{Z_N(\beta)}{\E Z_N(\beta)}.
\]
Thus $W_N=\log(1+X_N)$, and Theorem~\ref{thm:partition-clt} gives
\[
  N^aX_N
  \Longrightarrow
  \sqrt{\frac{q!}{2}}\,\beta\ell_q'(\beta)\,\xi,
  \qquad
  \delta_N^2:=\E X_N^2\asymp N^{-(q-2)}.
\]
We claim that
\begin{equation}\label{eq:log-centering-negligible}
  \E W_N=o(N^{-(q/2-1)}).
\end{equation}
To prove this, we first establish an $N$-independent Lipschitz bound for
$g\mapsto W_N(g)$.  Differentiation gives
\[
  \partial_{g_I}W_N
  =-\beta\sigma_N\langle\psi_I\rangle_{\beta,g}.
\]
For each $(q-1)$-set $S$, the operators
$(\psi_{S\cup\{i\}})_{i\notin S}$ are pairwise anticommuting Hermitian
involutions.  If $A_1,\ldots,A_m$ are any such operators and
$a_j=\langle A_j\rangle$ in a state, set $B=\sum_j a_jA_j$ and
$s=\sum_j a_j^2$.  Pairwise anticommutation gives
\[
  B^2=s\I,
  \qquad
  s=\langle B\rangle\leq\norm{B}=\sqrt{s},
\]
and hence $s\leq1$.  Applying this observation for each $S$ and noting that
every $I\in\calI_N$ is counted exactly $q$ times yields
\[
  \sum_{I\in\calI_N}|\langle\psi_I\rangle_{\beta,g}|^2
  \leq \frac1q\binom N{q-1}.
\]
Consequently,
\[
  \sum_{I\in\calI_N}|\partial_{g_I}W_N|^2
  \leq
  \frac{\beta^2\sigma_N^2}{q}\binom N{q-1}
  \leq\frac{\beta^2}{q}.
\]
Writing $\mu_N=\E W_N$, the Gaussian concentration inequality therefore
implies, for every $\lambda\in\mathbb R$,
\begin{equation}\label{eq:WN-subgaussian}
  \E e^{\lambda(W_N-\mu_N)}
  \leq \exp\left(\frac{\lambda^2\beta^2}{2q}\right).
\end{equation}
Jensen's inequality and the identity $\E e^{W_N}=1$ give
\[
  -\frac{\beta^2}{2q}\leq \E W_N\leq0.
\]
Indeed, the upper bound follows from Jensen, while
\eqref{eq:WN-subgaussian} with $\lambda=1$ gives the lower bound.

Let $K_\beta=\beta^2/(2q)$.  Since $e^{W_N}-1=X_N$, Chebyshev's inequality
gives, for $t\geq1$,
\[
  \mathbb P(W_N<-t)
  \leq \frac{\delta_N^2}{(1-e^{-t})^2}
  \leq C\delta_N^2.
\]
On the other hand, $\mu_N\geq-K_\beta$, so
\eqref{eq:WN-subgaussian} and the Chernoff bound imply
\[
  \mathbb P(W_N<-t)
  \leq
  \exp\left(-\frac{q(t-K_\beta)^2}{2\beta^2}\right),
  \qquad t>K_\beta.
\]
Set
\[
  T_N=\max\left\{
    1,\,
    K_\beta+\frac{2|\beta|}{\sqrt q}
    \sqrt{\log(e/\delta_N)}
  \right\}.
\]
Indeed, the layer-cake formula gives
\[
  \E\bigl[|W_N|\,\mathbf1_{\{W_N<-1\}}\bigr]
  =
  \mathbb P(W_N<-1)+
  \int_1^\infty\mathbb P(W_N<-t)\,\mathrm dt.
\]
The interval $[1,T_N]$ contributes at most $C\delta_N^2T_N$, while
integration of the Gaussian bound over $[T_N,\infty)$ contributes at most
$C\delta_N^2$.  Consequently,
\begin{equation}\label{eq:WN-negative-tail}
  \E\bigl[|W_N|\,\mathbf1_{\{W_N<-1\}}\bigr]
  \leq C\delta_N^2\sqrt{\log(e/\delta_N)}.
\end{equation}

Now set $h(w)=e^w-1-w$.  There is an absolute constant $C$ such that
\[
  h(w)\leq C(e^w-1)^2\quad (w\geq-1),
  \qquad
  h(w)\leq |w|\quad (w<-1).
\]
Since $\E e^{W_N}=1$, \eqref{eq:WN-negative-tail} yields
\[
  0\leq-\mu_N=\E h(W_N)
  \leq C\delta_N^2\sqrt{\log(e/\delta_N)}
  =O\!\left(N^{2-q}\sqrt{\log N}\right)
  =o(N^{-(q/2-1)}).
\]
This proves \eqref{eq:log-centering-negligible}.

It remains to compare the random parts.  The sequence $N^aX_N$ is tight and
$X_N\to0$ in probability.  On the event $\{|X_N|\leq1/2\}$, the Taylor
bound $|\log(1+x)-x|\leq Cx^2$ gives
\[
  N^a|W_N-X_N|
  \leq C\bigl(N^a|X_N|\bigr)|X_N|.
\]
The right-hand side converges to zero in probability, while
$\mathbb P(|X_N|>1/2)\to0$.  Therefore
\[
  N^a\bigl(W_N-X_N\bigr)\longrightarrow0
  \qquad\text{in probability}.
\]
Together with \eqref{eq:log-centering-negligible}, this shows that
\[
  N^a\bigl((W_N-\E W_N)-X_N\bigr)
  \longrightarrow0
  \qquad\text{in probability}.
\]
Finally,
\[
  \log Z_N-\E\log Z_N=W_N-\E W_N,
\]
so the claimed convergence follows from
\eqref{eq:scaled-relative-clt}.
The assertion for $\mathcal G_N$ follows by multiplication by
$-\beta^{-1}$.
\end{proof}
\begin{remark}
Since $\ell_q'(\beta)=\beta/q+O(\beta^3)$ at the origin, the nondegeneracy
condition in Theorem~\ref{thm:partition-clt} holds for every sufficiently
small real $\beta\neq0$.
\end{remark}


\begin{thebibliography}{99}

\bibitem{SachdevYe1993}
S.~Sachdev and J.~Ye,
\newblock Gapless spin-fluid ground state in a random quantum Heisenberg
magnet,
\newblock \emph{Phys. Rev. Lett.} 70 (1993), 3339--3342.
\newblock \href{https://doi.org/10.1103/PhysRevLett.70.3339}
{doi:10.1103/PhysRevLett.70.3339}.

\bibitem{MaldacenaStanford2016}
J.~Maldacena and D.~Stanford,
\newblock Remarks on the Sachdev--Ye--Kitaev model,
\newblock \emph{Phys. Rev. D} 94 (2016), 106002.
\newblock \href{https://doi.org/10.1103/PhysRevD.94.106002}
{doi:10.1103/PhysRevD.94.106002}.

\bibitem{FTW2018}
R.~Feng, G.~Tian, and D.~Wei,
\newblock Spectrum of SYK model II: Central limit theorem,
\newblock \emph{Random Matrices Theory Appl.} 10 (2021), no.~4, 2150037.
\newblock \href{https://doi.org/10.1142/S2010326321500374}
{doi:10.1142/S2010326321500374}.

\bibitem{Zlokapa2026}
A.~Zlokapa,
\newblock A rigorous quasipolynomial-time classical algorithm for SYK thermal
expectations,
\newblock \href{https://arxiv.org/abs/2604.21089}
{arXiv:2604.21089v1}, 2026.

\bibitem{GPSZ2026}
D.~Gamarnik, F.~Pernice, A.~Schmidhuber, and A.~Zlokapa,
\newblock The free energy limit of the SYK model at high temperature,
\newblock \href{https://arxiv.org/abs/2605.02768}
{arXiv:2605.02768v1}, 2026.

\bibitem{AGK2025}
E.~R. Anschuetz, D.~Gamarnik, and B.~T. Kiani,
\newblock Bounds on the ground state energy of quantum $p$-spin Hamiltonians,
\newblock \emph{Commun. Math. Phys.} 406 (2025), article 232.
\newblock \href{https://doi.org/10.1007/s00220-025-05412-4}
{doi:10.1007/s00220-025-05412-4}.

\end{thebibliography}

\end{document}
